\begin{proposition}[Known with derivation: identification with the opened joint Earle--Marden original]\label{A1-Earle-Marden-original}
The joint manuscript of Earle and Marden,
\emph{Holomorphic Plumbing Coordinates on Teichm\"uller and Compactified
Moduli Space}, dated \emph{Draft 8/24/11}, gives the following identification.
It is an original author manuscript, recovered from the public archive of
Marden's Minnesota papers directory, rather than the distinct published
2012 chapter. The relevant locators are \S4.1.1--4.1.2, pp.~39--40;
\S4.5, p.~50; and \S4.8.1--4.8.2, equations~(4.8.1)--(4.8.2),
p.~57 \cite{EarleMarden2011DraftA567}.

Normalize the peripheral translation to $A(\zeta)=\zeta+1$.
Suppose the two disjoint puncture neighborhoods are the quotients of
$\operatorname{Im}\zeta>h$ and $\operatorname{Im}\eta<-h$, where $h>0$.
For $\tau=a+2si$, $s>h$, the identification
$\zeta=\eta+\tau$ becomes classical plumbing
\[
 z=e^{2\pi i\zeta},\qquad w=e^{-2\pi i\eta},\qquad
 zw=t,\qquad t=e^{2\pi i\tau}.
\]
Writing $\rho=e^{-2\pi h}$ and $r=e^{-2\pi s}$ gives the exact bounds
\[
 0<r<\rho<1,\qquad |t|=r^2<\rho^2=e^{-4\pi h}.
\]
The sewn collar is the annulus
\[
 \frac{|t|}{\rho}<|z|<\rho,\qquad
 \operatorname{mod}=\frac{1}{2\pi}\log\frac{\rho^2}{|t|}
                   =\operatorname{Im}\tau-2h.
\]
On its middle circle the sewing is
$re^{i\theta}\mapsto re^{i(\phi-\theta)}$, where
$t=|t|e^{i\phi}$. More generally, after rescaling by positive
$0<r_+,r_-<\rho$ satisfying $r_+r_-=|t|$, put $u=z/r_+$ and $v=w/r_-$.
Then the same sewing is
\[
 uv=e^{i\phi},\qquad
 t\,\partial_z\otimes\partial_w
   =e^{i\phi}\,\partial_u\otimes\partial_v.
\]
Thus the toric boundary welding in Proposition~\ref{A1-welding-plumbing}
is exactly this plumbing operation when its full coordinate germs are
retained. The normal tensor is its first-order invariant.
\end{proposition}
\begin{proof}
The exponential maps are invariant under $A$, have nonzero derivatives,
and give biholomorphisms from the two half-plane quotients to the
punctured disc of radius $\rho$. Their signs follow directly from
$|e^{2\pi i\zeta}|=e^{-2\pi\operatorname{Im}\zeta}$.
On the identified points $\zeta-\eta=\tau$, so multiplication gives
$zw=e^{2\pi i\tau}$. The two cutting lines
$\operatorname{Im}\zeta=s$ and $\operatorname{Im}\eta=-s$ map to
the circles of radius $r$. Since $|t|=e^{-2\pi\operatorname{Im}\tau}
=e^{-4\pi s}$, these radii and the collar inequalities are exact.

Deleting the smaller discs and identifying their boundary circles by
$w=t/z$ gives the same analytic surface as identifying the overlapping
annuli $|t|/\rho<|z|<\rho$ and $|t|/\rho<|w|<\rho$.
Indeed the transition is holomorphic with nonzero derivative and
continues across the seam; the two half-annuli form this one annulus.
Its modulus is the displayed logarithmic ratio. The phase formula and
rescaled sewing equation follow by substitution. Finally,
$\partial_u=r_+\partial_z$ and
$\partial_v=r_-\partial_w$, proving the tensor identity.
Under any further changes $z'=bz+O(z^2)$, $w'=cw+O(w^2)$,
the first-order smoothing scalar is multiplied by $bc$ while the
tangent basis is multiplied by $(bc)^{-1}$. Consequently the tensor
is independent of those coordinate descriptions. No assertion about
the higher coordinate jets follows from this first-order invariance.
\end{proof}

\paragraph{What was checked in the original.}
The source's unambiguous classical construction is $zw=s$ in
\S4.8.1, equation~(4.8.1), followed by cutting at radius $\sqrt{|s|}$
and sewing by $w=s/z$ in equation~(4.8.2). Its \S6.1.1, p.~67,
also writes the annular bundle explicitly as
$p(z,w)=zw$ on two punctured unit discs.
The printed p.~50 image was inspected: its upper/lower half-plane
exponential signs and its displayed angular expression are inconsistent
with its own disc and sewing conventions. The signs and angle in the
proposition are obtained by the direct calculation above; they are
not silently copied from those draft misprints. The exact
$e^{-4\pi h}$ bound follows from the exponential relation; the printed
$e^{-4h}$ bound is merely weaker.

For comparison, the annular-cover argument in \S4.1(III), p.~42,
equation~(4.1.1), places its cyclic cover between strips of widths
$\operatorname{Im}\tau-2h$ and $\operatorname{Im}\tau+2h$.
Whenever those stated strip inclusions hold, monotonicity of annulus
modulus and $\ell=\pi/\operatorname{mod}$ for the core of a complete
hyperbolic annulus give
\[
 \frac{\pi}{\operatorname{Im}\tau+2h}
 \leq\ell\leq
 \frac{\pi}{\operatorname{Im}\tau-2h}.
\]
In particular, for this fixed strip bound,
$\ell\log(1/|t|)\longrightarrow2\pi^2$ as
$\operatorname{Im}\tau\longrightarrow\infty$, the source's
Proposition~4.5.1, equation~(4.5.1), pp.~50--51.
The coordinate identification itself uses only the explicit local
construction proved above.

\paragraph{Exact source status and scope.}
This directly opened joint original discharges the requested
identification with Earle--Marden's plumbing construction.
Its global statements, including Theorem~4.6.1, p.~51, have not been
used as proof inputs here. They are distinct from the directly verified
local operation, and the draft itself warns in \S4.8.3, p.~58, that
arbitrary local coordinates do not give global plumbing coordinates.
The separate published chapter \emph{Holomorphic plumbing coordinates},
\emph{Contemporary Mathematics}~575 (2012), pp.~41--52, remains
unopened in this investigation. The draft is neither identified with
that chapter nor assigned its publication status.
